% TOP_PROBLEM: {"id": 2400002, "title": "Fuglede's conjecture in dimensions one and two", "category_id": 8, "category": {"id": 8, "name": "analysis", "display_name": "Analysis", "description": "Limits, continuity, calculus, and function theory.", "slug": "analysis", "order_index": 8, "created_at": "2026-07-31T15:26:25.671Z"}, "status": "partially_solved"}
% FIRST_POSTED: null
% ENTRY_KIND: statement_only
% Imported from: batch08.tex; UnsolvedMath ID 2400002
% Compile twice: lualatex 2400002.tex
\documentclass[11pt]{article}
\usepackage{fontspec}
\setmainfont{Latin Modern Roman}
\usepackage{amsmath,amssymb,mathtools,unicode-math}
\setmathfont{Latin Modern Math}
\usepackage{xurl,hyperref}
\hypersetup{hidelinks,pdftitle={UnsolvedMath Problem 2400002}}
\newcommand{\sref}[2]{\hyperlink{source:#1:#2}{[S#2]}}
\newcommand{\cataloguescope}[1]{\paragraph{Mathematical question.}#1}
\begin{document}
% UnsolvedMath ID 2400002; source catalogue code AMR-023-0002
\setcounter{section}{2400001}
\section{Fuglede's conjecture in dimensions one and two}\label{2400002}
{\small\sffamily UnsolvedMath ID 2400002 · Analysis}
\subsection{Definitions and mathematical statement}

\paragraph{Shared notation.}$\mathbb N=\{1,2,\ldots\}$, and $\mathbb Z,\mathbb Q,\mathbb R,\mathbb C$ denote the integers, rationals, reals and complex numbers. Any other conventions are specified locally.


Let \(d\in\{1,2\}\) and let \(\Omega\subset\mathbb R^d\) be a Lebesgue-measurable set with \(0<|\Omega|<\infty\), without any convexity assumption. Call \(\Omega\) spectral if there is a countable set \(\Lambda\subset\mathbb R^d\) for which
\[
 \bigl\{|\Omega|^{-1/2}e^{2\pi i\langle\lambda,x\rangle}:\lambda\in\Lambda\bigr\}
\]
is an orthonormal basis of \(L^2(\Omega)\). Call \(\Omega\) a translational tile if there is a countable set \(T\subset\mathbb R^d\) such that
\[\sum_{t\in T}\mathbf1_\Omega(x-t)=1\quad\text{for Lebesgue-almost every }x\in\mathbb R^d.\]
Here \(|\Omega|\) denotes Lebesgue measure, \(\langle\cdot,\cdot\rangle\) the Euclidean inner product, and \(\mathbf1_\Omega\) the indicator function. These definitions concern equivalence up to null sets.
\paragraph{Statement recorded in the supplied source.}
In dimensions one and two, is every spectral set a translational tile, and is every translational tile spectral? The question concerns the unrestricted, possibly nonconvex, version of Fuglede's conjecture. \sref{2400002}{2}
\subsection{Short English statement}
Does having an orthogonal exponential basis characterize exactly the measurable sets that tile the line or plane by translations, when convexity is not assumed?
\subsection{Sources}
\begin{description}
\item[\hypertarget{source:2400002:1}{S1}] ULAM.AI and the UnsolvedMath Contributors, UnsolvedMath: A Curated Collection of Open Mathematics Problems, record 2400002 (AMR-023-0002). CC BY 4.0. \url{https://huggingface.co/datasets/ulamai/UnsolvedMath}.
\item[\hypertarget{source:2400002:2}{S2}] Dorin Ervin Dutkay and Palle E. T. Jorgensen, The momentum operator on a union of intervals and the Fuglede conjecture (2023), arXiv:2301.11377, Abstract and Section 1, pp. 1–2; Section 6. \url{https://arxiv.org/pdf/2301.11377}.
\end{description}
\label{end:2400002}
\end{document}
